%% ════════════════════════════════════════════════════════════════════════════
%% APPENDIX: Policy-Aligned Inpainter --- SFT and BCO Training Pipeline
%% ════════════════════════════════════════════════════════════════════════════

\section{Policy-Aligned Inpainter: SFT and BCO Training Details}

This appendix documents the complete training pipeline for the policy-aligned
inpainter $\mathcal{A}_g$, covering: the Refusal SFT stage
(\S\ref{app:sft}), the alignment objective selection and the migration from KTO to BCO (\S\ref{app:algo_choice}--\ref{app:kto_to_bco}), the full BCO loss formulation and motivations (\S\ref{app:bco_full}), the
failure archaeology from all training runs (\S\ref{app:failure_modes}), and
the final validated training configuration (\S\ref{app:final_config}).

%% ────────────────────────────────────────────────────────────────────────────
\subsection{Alignment Objective Selection: Why KTO and BCO}
\label{app:algo_choice}
%% ────────────────────────────────────────────────────────────────────────────

\paragraph{Dataset constraints.}
Modern RLHF methods impose different data-preparation burdens.
DPO~\cite{wallace2023diffusion_dpo} requires \emph{paired} (chosen, rejected)
completions for every prompt; for images this demands either human preference
labelling or an expensive oracle model.

KTO~\cite{li2024diffusion_kto} and BCO~\cite{jung2024bco} both operate on
\emph{unpaired} binary-labelled data $(x, y)\in\{0,1\}$, where $y=1$
denotes a safe (desirable) completion and $y=0$ denotes a policy-violating
(undesirable) one.
This format is produced directly by the auditor, requiring no
additional human annotation.

\paragraph{Risk profile motivation.}
For a content-safety inpainter, the error costs are highly asymmetric:
a false negative (generating adversarial material undetected) is far more
costly than a false positive (over-suppressing a benign region).
KTO was initially chosen specifically for its loss-averse utility function
and its asymmetric treatment of gains and losses, properties that appeared
well-suited to this risk profile.

%% ────────────────────────────────────────────────────────────────────────────
\subsection{Refusal SFT - LoRA Fine-Tuning}
\label{app:sft}
%% ────────────────────────────────────────────────────────────────────────────

\paragraph{Motivation: the know-before-you-refuse principle.}
The LLM training paradigm provides a useful analogy: instruction tuning
precedes alignment because a model must first \emph{know how to perform a
task} before it can refuse it in a principled way.
Applying an alignment signal directly to a model with no prior concept of
``safe inpainting'' is under-constrained: the model has no learned attractor
basin for policy-compliant completions, so the alignment gradient has
nothing to guide it toward.

Formally, let $\pi_\theta$ be the inpainter policy and $\mathcal{R}$ the
target ``refusal manifold'' in latent space---the set of completions that
replace NSFW content with spatially coherent, policy-compliant alternatives.
An alignment objective such as BCO applies a gradient
$\nabla_\theta \mathcal{L}_\text{BCO}$ that rewards distance from the
policy-violating attractor.
However, if $\mathcal{R}$ is not already in the support of $\pi_\theta$,
the gradient has no preferred direction in which to move the policy; it
merely pushes it \emph{away} from the violation without specifying
\emph{where} to go.
SFT pre-seeds the policy with a collapsed distribution over $\mathcal{R}$,
giving the BCO gradient a meaningful signal basin to expand.

In our training logs, this manifests clearly: the BCO run starting from an
SFT-initialised checkpoint exhausts its SFT prior through approximately
step~500, then searches the manifold for an exploit.
The asymmetric auxiliary losses (reconstruction and identity, detailed
in \S\ref{app:bco_full}) constrain this search to a corridor that
eventually leads to the globally optimal solution---rendering coherent
clothing.
A BCO run starting from a vanilla SD-inpainting checkpoint does not exhibit
this three-phase structure; it collapses directly to the pixel-smudging
attractor with no recovery.

\paragraph{Base model and adapter.}
SFT training starts from \emph{runwayml/stable-diffusion-inpainting}, a
9-channel inpainting UNet conditioned on
$(z_t \in \mathbb{R}^{4}, m_\ell \in \mathbb{R}^{1}, \tilde{z}_0 \in
\mathbb{R}^{4})$, where $m_\ell$ is the latent-space mask and $\tilde{z}_0$
is the masked image latent.
We apply LoRA~\cite{hu2022lora} with rank $r{=}64$, $\alpha{=}64$
(effective scale $= \alpha/r = 1.0$, i.e.\ no implicit scaling)
targeting the following parameter groups:
\begin{itemize}[leftmargin=1.5em,itemsep=1pt,topsep=2pt]
  \item Cross- and self-attention:
        \emph{to\_q}, \emph{to\_k}, \emph{to\_v}, \emph{to\_out.0}
  \item Feed-forward layers in transformer blocks:
        \emph{ff.net.0.proj}, \emph{ff.net.2}
  \item Spatial projection layers: \emph{proj\_in}, \emph{proj\_out}
  \item ResNet convolutional layers (structural/spatial output):
        \emph{conv1}, \emph{conv2}
  \item UNet input convolution layer: \texttt{conv\_in}
\end{itemize}

The inclusion of \emph{conv\_in} is critical.
The 9-channel inpainting UNet differs from the standard 4-channel text-to-image
UNet in precisely this first layer: it processes the concatenated
$(z_t, m_\ell, \tilde{z}_0)$ tensor.
Leaving \emph{conv\_in} frozen during alignment causes the
mask-conditioning logic to resist policy signals, because the gradient
from the loss cannot propagate into the layer that interprets which pixels
to modify.
Similarly, the ResNet convolutional layers (\emph{conv1}, \emph{conv2})
handle the structural and spatial content of the output.
Refusal requires rendering new content (clothing, blurred faces)---a
structural operation---and attention-only LoRA lacks the representational
capacity for such changes.

Stage B-2 uses rank $r{=}128$, $\alpha{=}128$ for the alignment phase,
on the hypothesis that the earlier runs with smaller rank were exhausting
LoRA capacity and saturating prematurely.

\paragraph{Training data.}
The SFT dataset consists of approximately 2\,000 images (1\,000 nudity
violations, 1\,000 violence violations).
For each unsafe source image, the auditor generates a mask over the
policy-violating region.
The source image and mask are then passed to the inpainter with a
custom-generated safe prompt (derived by prompting a Qwen-2.5-VL captioner
followed by a LLaMA-8B rewriter) to produce a target ``accepted'' image.
Training pairs are therefore $(x_\text{unsafe}, m, x_\text{safe}, p_\text{unsafe})$:
the unsafe source, its violation mask, the human-approved safe target, and
the original adversarial prompt.

Training on the \emph{adversarial} prompt rather than the safe prompt is
deliberate: at inference, the inpainter receives the adversarial prompt
because the user's intent is unknown.
The model must learn the pixel-space signature of a safe completion
conditioned on an unsafe prompt---not relying on text-based cues to trigger
refusal.

\paragraph{Loss and schedule.}
The SFT loss is Min-SNR~\cite{hang2023minsnr} weighted MSE across the full
noise schedule $t \in [0, 1000)$:
\[
  \mathcal{L}_\text{SFT} = \mathbb{E}_{t,\epsilon}\!\left[
    w(t)\cdot
    \|\epsilon_\theta(\begin{aligned}z_t, t, m_\ell, \\ \tilde{z}_0, p_\text{unsafe} \end{aligned}) 
     - \epsilon\|^2
  \right],
  \label{eq:sft_loss}
\]
where $w(t) = \min(\text{SNR}(t), \gamma)/\text{SNR}(t)$ with $\gamma{=}5$.
Without Min-SNR reweighting, high-noise timesteps (large $t$) dominate the
gradient because the signal-to-noise ratio is low and the loss magnitude is
large.
This matters here because safe completions are visually coherent only at
\emph{low} noise levels ($t \ll 500$); high-noise gradients contribute
primarily noise-level statistics and not the structural pattern of
clothing or background fill.
Min-SNR redistributes gradient mass toward the low-noise regime where the
visual content of the refusal is actually resolved.

Formally, the Min-SNR weight is derived from the signal-to-noise ratio
$\text{SNR}(t) = \bar\alpha_t / (1 - \bar\alpha_t)$, where
$\bar\alpha_t = \prod_{s=1}^t \alpha_s$ is the cumulative product of the
noise schedule.
The clamp at $\gamma$ prevents the weight from collapsing to zero at very
low $t$ (where SNR is enormous), maintaining gradient flow throughout the
schedule.

A small noise offset of $0.05$ is added:
$\epsilon' = \epsilon + 0.05 \cdot \epsilon_0$ where $\epsilon_0 \sim
\mathcal{N}(0, I_{C \times 1 \times 1})$ is broadcast.
This offset helps with dark and saturated regions---common in clothing
inpainting---by preventing the VAE latent distribution from collapsing into
degenerate attractors near the boundary of its quantization grid.

After convergence (loss plateau below $\sim$0.07 over 8 epochs), the LoRA
weights are merged into the base checkpoint via \texttt{merge\_and\_unload()},
producing a single 9-channel UNet used as the Stage B-2 initialisation.

\paragraph{How SFT initialises BCO: a geometric argument.}
Let $\mathcal{M}_\text{img}$ denote the natural image manifold in latent
space and $\mathcal{M}_\text{safe} \subset \mathcal{M}_\text{img}$ the
submanifold of safe completions inside the mask.
Pre-SFT, $\pi_\theta$ (the base inpainting model) has learned a strong
prior toward $\mathcal{M}_\text{img}$ conditioned on the input context, but
has never encountered a training signal that distinguishes
$\mathcal{M}_\text{safe}$ from the rest of the manifold.

SFT provides approximately 2\,000 examples of $(x_\text{unsafe}, m,
x_\text{safe})$ triples.
The MSE loss \eqref{eq:sft_loss} carves a narrow basin in the loss
landscape centred on $\mathcal{M}_\text{safe}$ for inputs of the form
$(z_t, m_\ell, \tilde{z}_0, p_\text{unsafe})$.
After SFT, $\pi_\theta$ has a collapsed mode at clothing/covering
completions for masked unsafe inputs.

BCO then operates on this initialised policy.
The BCO gradient moves the policy such that:
(a) on safe inputs, the policy stays near or improves upon the reference;
(b) on unsafe inputs, the policy deliberately degrades relative to the
reference (moves further from the policy-violating content).
Because the SFT basin already exists, the BCO search for ``worse on
unsafe'' lands directly on $\mathcal{M}_\text{safe}$ rather than on
degenerate off-manifold solutions.
This explains the three-phase dynamics observed in BCO training: the SFT prior
is exploited first (Phase 1), the degenerate attractor is blocked by
auxiliary losses (Phase 2), and the SFT basin is rediscovered as the
unique viable optimum (Phase 3).

%% ────────────────────────────────────────────────────────────────────────────
\subsection{KTO: Formulation, Experiments, and Failure Modes}
\label{app:kto_to_bco}
%% ────────────────────────────────────────────────────────────────────────────

We now document the KTO implementation and the concrete failure modes that motivated the switch to BCO.

\subsubsection{KTO Formulation}

The KTO objective~\cite{ethayarajh2024ktomodelalignmentprospect} for diffusion models \cite{li2024diffusion_kto} adapts the
Kahneman--Tversky prospect-theoretic value function to the noise-prediction setting.
Our implementation computes a ``g-term'' as:
\begin{equation}
  g_i = \ell^\theta_i - \ell^\text{ref}_i, \notag
  \label{eq:kto_gterm}
\end{equation}
where $\ell^\bullet_i$ is the masked MSE of model $\bullet$ against the
clean latent $z_0$ inside the mask:
\begin{equation}
  \ell^\bullet_i
  = \frac{\sum_{b,c,h,w} (1 + \lambda_m m_\ell) \cdot
      (\hat{z}_0^\bullet - z_0)^2 \cdot m_\ell}
    {\sum m_\ell}. \notag
  \label{eq:masked_mse}
\end{equation}
A positive $g_i$ means the trainable model has drifted further from $z_0$
than the reference inside the mask---\emph{worse} alignment with the safe target.

The signal is centered using an EMA over safe samples:
\[
  \delta_\text{ema} \leftarrow 0.99\,\delta_\text{ema}
    + 0.01\,\bar{g}|_{\text{safe}},
  \quad \label{eq:kto_ema}
\]
\[
  \delta_\text{ema} \in [-0.02, 0.05].
\]
The centered signal drives a sigmoid-based loss:
\begin{equation}
  h_i = \sigma\!\left(\text{sgn}_i \cdot \beta \cdot (g_i - \delta_\text{ema})\right),
  \quad \notag
  \end{equation}
  \begin{equation}
  \mathcal{L}_\text{KTO} = \mathbb{E}_i\!\left[w_y \cdot (1 - h_i)\right], \notag
  \label{eq:kto_loss}
\end{equation}
where $\text{sgn}_i = +1$ for safe samples and $-1$ for unsafe,
$\beta = 7$, and $w_y \in \{1.0, 5.0, 12.0\}$ for safe, nudity, and
violence classes respectively.

\subsubsection{KTO Experiment Progression}

\paragraph{Vanilla KTO.}
The first KTO implementation used a static reference model and the
$\epsilon$-space MSE as a proxy for the probability of generating
policy-violating content.
The model quickly found a degenerate attractor: generating uniform blue
patches inside the masked region.
This trivially reduced the noise-prediction MSE difference relative to the
reference (whose MSE was also high at the same patches) while producing
visually incoherent output.
The model had hacked the reward by leaving the image manifold entirely.

\paragraph{Dynamic reference.}
To prevent manifold departure, we tried updating the reference model every
$N$ gradient steps.
This accelerated instability: accumulated errors in the trainable model
were passed to the reference model, which then guided training toward its
own failures in a positive-feedback loop.
Image quality outside the mask degraded correspondingly.

\paragraph{Reconstruction and identity losses.}
To anchor the model to the image manifold, we introduced two auxiliary
losses (described fully in \S\ref{app:bco_full}):
a reconstruction loss penalising noise-prediction error outside the mask,
and an identity guardrail penalising $\hat{z}_0$-space divergence from the
reference.
These improved manifold stability but introduced a new failure mode: the
model learned to smudge hands, feet, and peripheral body parts inside the
mask---a pixel-destruction strategy that achieved low KTO/BCO loss while
staying near the identity boundary.

\paragraph{Modified loss formulation.}
We observed that the reference model did not always generate a
policy-violating inpaint, especially for prompts that did not explicitly
mention NSFW content (Qwen-2.5-VL captions tend to be descriptive rather
than explicit).
We reframed the problem as alignment \emph{and} safety-quality improvement,
computing the g-term as the difference of MSEs against ground-truth noise
rather than against the reference.
Results were similar.

\paragraph{Prompt dropout.}
Observing that the cross-attention layers were overweighted (higher
activation delta than before the manifold dip and after recovery), we
suspected the cross-attention was learning to key on NSFW token patterns
rather than visual features.
We introduced 30\% prompt dropout per batch, zeroing the text conditioning
to force spatial grounding of the refusal signal.
This engineering choice is retained in the BCO configuration.

\subsubsection{Why KTO Failed: Theoretical Analysis}

\paragraph{Failure Mode 1: $\epsilon$-space stochasticity.}
The KTO g-term \eqref{eq:kto_gterm} is computed from $\hat{z}_0$
predictions, but the $\hat{z}_0$ reconstruction is itself a function of
the noise $\epsilon$ sampled at each step.
At high timesteps $t \sim \mathcal{U}[500, 1000)$, the standard deviation
of $\epsilon$ is $\sim\!\mathcal{O}(1)$, and the resulting variance in the
masked MSE across samples from the same image at the same $t$ is
$\mathcal{O}(10^{-1})$.
The preference gap between safe and unsafe images in our dataset is
$\mathcal{O}(10^{-2})$.
Consequently, the signal-to-noise ratio of the preference gradient is less
than 1; the loss tracks noise-schedule variance rather than safety
information.
This produces the characteristic flat
$h_\text{safe} \approx h_\text{unsafe} \approx 0.500$ plateau observed in
early KTO runs.

\paragraph{Failure Mode 2: Asymmetric risk cannot be encoded cleanly.}
KTO uses $(1 - \sigma(x))$ as its loss function, which is equivalent to
$\log \sigma(-x)$ up to sign.
The curvature of this function is symmetric in $x$: moving the policy away
from a violation provides the same gradient magnitude as moving it toward
a safe completion.
Our risk profile is strongly asymmetric: false negatives on unsafe content
(NSFW output reaches the user) are far more costly than false positives
(over-suppression of benign content).
KTO has no mechanism to encode this asymmetry beyond the class weights
$w_y$, which scale the loss magnitude but not the gradient shape.

Additionally, KTO's EMA baseline is computed from safe samples only
(\eqref{eq:kto_ema}), then subtracted from both safe \emph{and} unsafe
g-terms.
This conflates the distributional statistics of two distinct populations:
safe images (where $g \approx 0$ under a well-behaved policy) and unsafe
images (where $g \gg 0$ is desirable).
Subtracting a safe-derived baseline from the unsafe signal systematically
miscalibrates the centering.

\paragraph{Failure Mode 4: Identity loss bleeds globally.}
The KTO implementation computes the identity gap as a single global scalar:
\begin{equation}
  \ell_\text{id} = \big(\|{\hat{z}_0^\theta - \hat{z}_0^\text{ref}}\|^2 - 0.02\big)_+^2, \notag
\end{equation}
aggregated over all samples.
A single severely drifted sample inflates this scalar, raising the penalty
for all other samples and causing the optimizer to over-regularize the
well-behaved ones.
The resulting signal is poorly calibrated and impedes the policy's ability
to diverge from the reference on unsafe samples---which is precisely the
desired behavior.

%% ────────────────────────────────────────────────────────────────────────────
\subsection{BCO Loss: Full Formulation with Proofs and Motivations}
\label{app:bco_full}
%% ────────────────────────────────────────────────────────────────────────────

BCO~\cite{jung2024bco} reformulates the alignment objective as binary
classification: a sample is ``desirable'' ($y=1$) if it should be
encouraged and ``undesirable'' ($y=0$) if it should be suppressed.
The following sections derive each component of the BCO loss from first
principles.

\subsubsection{Why $\hat{z}_0$, Not $\epsilon$}

Standard diffusion training minimises noise-prediction MSE
$\|\epsilon_\theta - \epsilon\|^2$.
At timestep $t$, the magnitude of $\epsilon$ is $\sim\!\mathcal{O}(1)$, so
small changes in $\epsilon_\theta$ produce large, timestep-dependent changes
in the MSE.
A reward computed from $\|\epsilon_\theta - \epsilon\|^2$ is therefore
dominated by $t$, making per-class weighting meaningless---a sample at
$t=900$ produces 10--100$\times$ the loss magnitude of the same sample at
$t=100$.

Instead, both UNets predict the clean latent:
\begin{equation}
  \hat{z}_0^\bullet
  = \frac{z_t - \sqrt{1-\bar\alpha_t}\,\epsilon_\bullet}{\sqrt{\bar\alpha_t}},
  \quad \bullet \in \{\theta, \text{ref}\}. \notag
  \label{eq:z0hat}
\end{equation}
The denominator $\sqrt{\bar\alpha_t}$ ranges from approximately $1.0$
(at $t=0$) to $\sim$0.01 (at $t=999$).
While $\hat{z}_0$ predictions are noisier at high $t$, they are calibrated
to the same latent space regardless of $t$: $\hat{z}_0 \in
[-3\sigma_{z_0}, +3\sigma_{z_0}]$ with bounded variance.
This invariance is the key property: per-class weights now operate on a
reward signal with consistent scale, making nudity weight $5.0\times$ and
violence weight $12.0\times$ interpretable as actual preference ratios.

\subsubsection{Masked MSE Reward}

The reward quantifies how much better the trainable policy reconstructs
the safe target inside the mask, relative to the frozen reference:
\begin{align}
  w &= 1 + \lambda_m \cdot m_\ell, \quad \lambda_m = 0.5, \notag
  \\
  \ell_\bullet &= \frac{\sum_{b,c,h,w} w \cdot (\hat{z}_0^\bullet - z_0)^2 \cdot m_\ell}
                       {\sum m_\ell}, \notag
                       \label{eq:masked_mse_bco} \\
  g_i &= \ell^\text{ref}_i - \ell^\theta_i. \notag 
\end{align}
Note the sign convention: $g_i > 0$ means the trainable model reconstructs
the masked region \emph{better} than the reference (closer to $z_0$).
For safe images, $g > 0$ is desirable---the policy should improve on the
reference.
For unsafe images, $g < 0$ is desirable---the policy should produce
reconstructions \emph{further} from the policy-violating $z_0$.

The mask weight $\lambda_m = 0.5$ applies $1.5\times$ gradient emphasis
inside the adversarial region without completely ignoring the surrounding
context.
Ignoring the unmasked region entirely ($\lambda_m \to \infty$, i.e.\ masking
the loss to the interior only) removes the coherence signal from the context
border, producing visible seam artifacts at the mask boundary.

\paragraph{Hinge cap on unsafe reward.}
For unsafe samples ($y=0$), the reward is capped at:
\begin{equation}
  g_i \leftarrow \min\!\Big(g_i,\;
    1.5 \cdot |\bar{g}_\text{unsafe}|_\text{detach}\Big), \notag
  \label{eq:hinge_cap}
\end{equation}
where $\bar{g}_\text{unsafe}$ is the mean reward over unsafe samples in
the current batch.
Without this cap, the policy can exploit the mask by uniformly suppressing
\emph{all} content inside it---a trivial maximisation of $-g$ that
achieves a very negative reward (far from $z_0$) while producing visually
incoherent output.
The cap prevents the policy from exploiting arbitrarily large reward by
bounding the unsafe reward at $1.5\times$ its own class mean,
forcing it to find solutions that are both safe-suppressing \emph{and}
spatially coherent.

\subsubsection{Reward-Shift EMA and BCO Loss}

To center the reward at the decision boundary between safe and unsafe
classes, we maintain an exponential moving average:
\begin{align}
  \delta_\text{raw} &= \tfrac{1}{2}
    \Big(\mathbb{E}[g | y{=}1] + \mathbb{E}[g | y{=}0]\Big), \notag
    \label{eq:delta_raw} 
\end{align}
\[
  \delta_\text{ema} \leftarrow
    0.999\,\delta_\text{ema} + 0.001\,\delta_\text{raw}, \label{eq:delta_ema_1}
\]
\[
  \quad \delta_\text{ema} \in [-0.03, +0.03]. \label{eq:delta_ema_2}
\]

\paragraph{Why the symmetric midpoint?}
The shifted reward $r_i = g_i - \delta_\text{ema}$ should be positive for
safe samples (the policy is better than baseline) and negative for unsafe
samples (the policy is worse than baseline).
The optimal centering point is the midpoint of the two class means, which is
exactly $\delta_\text{raw}$.
KTO's centering (\eqref{eq:kto_ema}) uses a safe-only EMA, which
systematically underestimates the true decision boundary when $g$ differs
across classes, miscalibrating the loss.

The EMA clamp is applied \emph{in-place} to the stored EMA value before any downstream use:
The $[-0.03, +0.03]$ bound is tighter and symmetric compared to KTO's $[-0.02, +0.05]$, preventing runaway baseline drift in either direction.

\paragraph{BCO loss: softplus vs.\ sigmoid.}
Given $r_i = g_i - \delta_\text{ema}$ and label sign
$\text{sgn}_i = 2y_i - 1 \in \{-1, +1\}$, the per-sample BCO loss is:
\begin{equation}
\begin{aligned}
  \ell^\text{BCO}_i = \mathrm{softplus}(-\,\text{sgn}_i \cdot \beta \cdot r_i)
  \\= \log(1 + \exp(-\,\text{sgn}_i \cdot \beta \cdot r_i)). \notag
  \label{eq:bco_loss}
\end{aligned}
\end{equation}

\paragraph{Why softplus instead of KTO's $(1 - \sigma)$ loss?}
Note that $\log \sigma(x) = -\text{softplus}(-x)$.
The KTO per-sample loss $w_y(1 - h_i)$ with $h_i = \sigma(\cdot)$ is
related to the negative log-likelihood $-\log \sigma(\cdot)$ but has a
different gradient profile.

\begin{enumerate}[leftmargin=1.5em,itemsep=2pt]
  \item \textbf{Gradient saturation.}
  $\partial(1-\sigma(x))/\partial x = -\sigma(x)(1-\sigma(x))$, which
  collapses to zero as $x \to \pm\infty$.
  $\partial\,\mathrm{softplus}(-x)/\partial x = -\sigma(-x) = -(1-\sigma(x))$,
  which goes to $-1$ as $x \to +\infty$ and to $0$ as $x \to -\infty$.
  Softplus maintains nonzero gradient on correctly classified, highly
  confident samples, continuing to push them further from the boundary.
  KTO's $(1-\sigma)$ loss provides a vanishing gradient once $h$ is large,
  causing training to stall for well-separated samples while allowing
  hard cases to drift.

  \item \textbf{Probabilistic interpretation.}
  Softplus is the numerically stable form of the binary cross-entropy
  loss $-\log \sigma(\cdot)$, which is the maximum-likelihood loss for a
  logistic classifier.
  KTO's $(1-\sigma)$ loss lacks this clean probabilistic interpretation
  and is less aligned with standard RL reward maximization intuition.

  \item \textbf{Numerical stability.}
  At $\beta=50$, the argument to the sigmoid can easily reach $|\beta r|
  \sim 50 \times 0.5 = 25$, causing floating-point overflow in
  $\exp(25) \approx 7 \times 10^{10}$.
  PyTorch's \texttt{F.softplus} handles this via the numerically stable
  formulation $\max(x, 0) + \log(1 + \exp(-|x|))$.
\end{enumerate}

The final BCO loss aggregates with per-class weights $w_y$:
\begin{equation}
\begin{aligned}
    \mathcal{L}_\text{BCO} = \mathbb{E}_i\!\left[
    w_y \cdot \mathrm{softplus}(-\,\text{sgn}_i \cdot \beta \cdot r_i)
  \right], \\
  \quad \beta = 50.
  \notag
  \label{eq:bco_loss_final} 
\end{aligned}
\end{equation}

\subsubsection{Identity Guardrail: From Linear Penalty to Quadratic Hinge}

\paragraph{The smudging problem.}
An early linear identity penalty $\mathcal{L}_\text{id} =
\|\hat{z}_0^\theta - \hat{z}_0^\text{ref}\|^2$ creates a continuous
tradeoff: the policy can achieve low BCO loss (suppressing unsafe content)
while staying near the reference (low identity loss) by \emph{smudging}
pixels---slightly blurring or washing out the masked region rather than
replacing it with coherent content.
Smudging produces $\|\hat{z}_0^\theta - \hat{z}_0^\text{ref}\|^2 \approx
0.01$--$0.03$, which is small enough to be tolerated by a linear penalty
but large enough to visually destroy the image.

\paragraph{The Quadratic Hinge.}
The fix is a threshold that allows free exploration below a safe-smudging
boundary and applies quadratic torque above it:
\begin{align}
  \ell_\text{id}(i) &=
    \big(\|\hat{z}_0^\theta(i) - \hat{z}_0^\text{ref}(i)\|^2
    - \kappa\big)_+^2, \quad \kappa = 0.02, 
    \notag
    \label{eq:id_hinge} \\
  \mathcal{L}_\text{id} &= \mathbb{E}_i\!\left[
    \bar{w}_i \cdot \ell_\text{id}(i)
  \right],
  \notag
\end{align}
where $\bar{w}_i = 30.0$ for safe samples and $\bar{w}_i = 5.0$ for unsafe .

The hinge at $\kappa = 0.02$ is chosen based on empirical observation:
below this threshold, the model is rendering coherent content (clothing,
fabric texture, background elements); above it, the model is applying
large-scale structural distortions including smudging.
Below $\kappa$: \emph{zero penalty}---the model may freely render
clothing, fabric, and background detail.
Above $\kappa$: \emph{quadratic torque}---the penalty grows as the
\emph{square} of the excess, making large distortions ($\ell_\text{id}
\gg \kappa$) prohibitively expensive.
A distortion of $\ell_\text{id} = 0.05$ incurs penalty $(0.05 - 0.02)^2
= 9 \times 10^{-4}$, while $\ell_\text{id} = 0.10$ incurs
$(0.10 - 0.02)^2 = 6.4 \times 10^{-3}$---a $7.1\times$ increase for a
$2\times$ increase in distortion.

\paragraph{Hinge applied per sample, not globally.}
BCO applies hinge \emph{per sample} before aggregating.
\begin{equation}
  \text{KTO: } \mathcal{L}_\text{id}
  = \Big(\mathbb{E}_i\!\left[\|\hat{z}_0^\theta - \hat{z}_0^\text{ref}\|^2\right]
    - 0.02\Big)_+^2.
    \notag
\end{equation}
With global aggregation, a single badly drifted sample inflates the mean,
raising the penalty for all other samples and causing the optimizer to
over-regularise the well-behaved ones.
Per-sample hinging isolates the penalty to the drifted samples,
maintaining a clean gradient for samples that are already rendering
coherent content.

\paragraph{Small weight for unsafe samples.}
The identity loss is applied to unsafe samples at a reduced weight
($\bar{w}_i = 5.0$ vs.\ $30.0$ for safe).
This is deliberate: for unsafe samples, the policy \emph{should} diverge
from the reference inside the mask (the reference produces NSFW content;
the policy should not).
However, the identity loss is applied across the full image ($\hat{z}_0$
is full-resolution), so it also anchors the \emph{unmasked} background.
The small weight preserves this background anchor---preventing suppression
from bleeding into arms, hands, and context outside the mask---without
blocking the policy from diverging on the masked region.

\subsubsection{Reconstruction Anchor}

The reconstruction loss anchors the unmasked background to the
ground-truth noise:
\begin{equation}
  \mathcal{L}_\text{recon}
  = \frac{\sum_i \bar{r}_i \cdot \|(\epsilon_\theta - \epsilon) \odot (1-m_\ell)\|^2_i}
    {\sum_i \bar{r}_i},
    \notag
  \label{eq:recon_loss}
\end{equation}
where $\bar{r}_i = 200.0$ for safe samples and $\bar{r}_i = 0.0$ for
unsafe.

\paragraph{Why noise space, not $\hat{z}_0$ space?}
The reconstruction loss operates in $\epsilon$-space rather than
$\hat{z}_0$-space, because it targets the \emph{unmasked} region where the
noise-schedule variance argument does not apply: at any timestep $t$, the
noise prediction on an unmasked region should match the ground-truth
noise $\epsilon$ regardless of $t$.
Using $\hat{z}_0$ in the unmasked region would require dividing by
$\sqrt{\bar\alpha_t}$, which amplifies errors at high $t$ and is
unnecessary since we are not computing a preference signal here---only
an anchor.

\paragraph{Why zero weight for unsafe?}
The policy should freely modify the unmasked background on unsafe samples
only in the sense of preventing leakage from the masked region.
Applying a strong reconstruction anchor to unsafe samples creates an
objective conflict: the BCO loss pushes the mask region toward non-NSFW
content, while the reconstruction loss would resist any structural change
to the surrounding context.
Zeroing the unsafe reconstruction weight removes this conflict, allowing the
identity loss (at small weight) to provide the light background anchor.

\paragraph{Normalized aggregation.}
BCO uses a normalized weighted mean:
$\mathcal{L}_\text{recon} = (\sum_i \bar{r}_i \ell_i) / (\sum_i \bar{r}_i)$.
Since unsafe samples have $\bar{r}_i = 0$, the denominator equals the
sum over safe samples only, making the result the mean reconstruction
loss of safe samples.
This is preferable to the KTO formulation of
$\mathcal{L}_\text{recon} = 200.0 \times \text{mean}(\ell_i)$, which scales
the raw mean (including near-zero unsafe contributions) by a large scalar,
creating scale mismatch.
Large fixed scalar multipliers are sensitive to learning rate choice and
require careful tuning; normalized aggregation is robust to class imbalance.

\subsubsection{Multi-Step DDIM Unrolling}

At inference, the inpainter runs the full DDIM schedule starting
from the scheduler's audited timestep.
Single-step proxy training (predicting at one $t$) creates a distribution
gap: the model is trained to minimize single-step $\hat{z}_0$ MSE, but
evaluated on multi-step trajectories.

Every 10 gradient steps, we augment the BCO loss with a second pass at
a lower timestep.
Given the current timestep $t$ and unrolling factor $K=10$:
\begin{align}
  \hat{z}_0^{(s)} &= \frac{z_{t_s} - \sqrt{1-\bar\alpha_{t_s}}\,
    \epsilon_\theta(z_{t_s}, t_s)}{\sqrt{\bar\alpha_{t_s}}},
    \notag\label{eq:ddim_z0} \\
  z_{t_{s+1}} &= \sqrt{\bar\alpha_{t_{s+1}}}\,\hat{z}_0^{(s)}
    + \sqrt{1-\bar\alpha_{t_{s+1}}}\,\epsilon_\theta(z_{t_s}, t_s),\notag
\end{align}
run for $K-1$ deterministic ($\eta=0$) DDIM steps under \texttt{no\_grad}
to produce $z_\text{mid}$.
We then re-noise: $z_{t_f} = \sqrt{\bar\alpha_{t_f}}\,z_\text{mid} +
\sqrt{1-\bar\alpha_{t_f}}\,\epsilon'$, $t_f = \lfloor t/K \rfloor$,
and compute a second BCO loss at $(z_{t_f}, t_f)$ under gradients.

Both passes contribute equally (weight $0.5$ each) to the final gradient:
\begin{equation}
  \mathcal{L}_\text{BCO}^\text{total}
  = 0.5\,\mathcal{L}_\text{BCO}^{(t)} + 0.5\,\mathcal{L}_\text{BCO}^{(t_f)}.
  \notag
\end{equation}
To prevent holding two full UNet activation graphs in VRAM simultaneously,
pass-1 backward is called before pass-2 forward.

\subsubsection{Full Objective}

\begin{equation}
  \mathcal{L} = 4\,\mathcal{L}_\text{BCO} + \mathcal{L}_\text{recon}
    + \mathcal{L}_\text{id}.
    \notag
  \label{eq:full_obj}
\end{equation}
The factor of $4$ on $\mathcal{L}_\text{BCO}$ was tuned empirically.
The softplus loss with $\beta=50$ has magnitude
$\sim\!\mathcal{O}(10^{-2}$--$10^{-1})$ at the decision boundary.
The reconstruction and identity losses have magnitude
$\sim\!\mathcal{O}(10^{-1}$--$10^{0})$ (noise-space MSE).
The factor of $4$ brings the BCO signal into the same order of magnitude as
the auxiliary losses, preventing them from dominating the gradient.

Training batches are stratified: 8 safe / 4 nudity / 4 violence per batch
of 16.
Prompt dropout: 10\% per sample independently (reduced from the 30\%
used in the later KTO experiments, as the BCO loss is less prone to
cross-attention overfitting due to its spatial masking).


% ─────────────────────────────────────────────────────────────────────────────
% FIGURE PLACEHOLDER — Model weight activations at steps 500 / 1000 / 2500
% ─────────────────────────────────────────────────────────────────────────────
\begin{figure*}[t]
  \centering
  \includegraphics[width=\linewidth]{latent-space-seperation.png}
  \caption{
    \textbf{The graph shows separation in the prompt trajectories in the latent space from high noise $t=700$ to low noise $t=5$ between red=`nudity', blue=`violence' and green=`safe' prompts for the policy-aware inpainter at step 0 and step 3000 of the BCO process for the attention module.}
  }
  \label{fig:activation_scatter}
\end{figure*}

%% ────────────────────────────────────────────────────────────────────────────
\subsection{Failure Mode Archaeology}
\label{app:failure_modes}
%% ────────────────────────────────────────────────────────────────────────────

Table~\ref{tab:failure_modes} documents the specific failure modes
discovered across the different training runs, the diagnostic evidence that
revealed them, and the architectural fix applied.

\begin{table*}[!ht]
\centering
\small
\setlength{\tabcolsep}{4pt}
\begin{tabular}{p{2.6cm} p{3.6cm} p{3.8cm} p{4.0cm}}
\toprule
\textbf{Failure Mode} & \textbf{Symptom} & \textbf{Root Cause}
  & \textbf{Fix Applied} \\
\midrule
Noise-space stochasticity (KTO Runs)
  & $h_\text{safe} \approx h_\text{unsafe} \approx 0.500$
    through step~500; no preference separation
  & Preference gap $\mathcal{O}(10^{-2})$ overwhelmed by
    high-$t$ noise variance $\mathcal{O}(10^{-1})$ in $\epsilon$-space;
    SNR of gradient $< 1$
  & Switch to BCO; reward computed on $\hat{z}_0$
    (bounded, schedule-invariant) \\
\addlinespace
Off-manifold blue patches (KTO, vanilla)
  & Model generates uniform blue fill inside mask;
    image quality deteriorates outside mask
  & Model hacks reward by leaving image manifold;
    noise-space MSE is trivially minimized by off-manifold output
  & Reconstruction + identity losses to anchor manifold \\
\addlinespace
Smudging / pixel destruction (KTO Runs and early BCO)
  & Unsafe content obscured by blur/smear rather than
    replaced; identity\_gap $\approx 0.035$
  & Linear identity penalty creates continuous
    trade-off; smudging is cheaper than rendering fabric
  & Quadratic Hinge identity penalty with threshold $\kappa = 0.02$ \\
\addlinespace
Mask-conditioning breakage
  & Model ignores mask; generates uniform texture across
    full image
  & Alignment gradients propagate into \texttt{conv\_in},
    corrupting 9-channel conditioning logic
  & Include \texttt{conv\_in} in LoRA adapter from SFT
    stage; $\mathcal{L}_\text{recon}$ anchors unmasked context \\
\addlinespace
Sigmoid saturation
  & \texttt{sigmoid\_sat\_pct} $> 90\%$ by step~1500;
    softplus gradient near floor; training stalls
  & $\beta{=}50$ with large reward gap pushes softplus
    into saturation; expected behavior, not a bug
  & Monitor $\Delta N$ for continued separation.
    Saturation = binary decision boundary learned.
    Reduce $\beta$ only if $\Delta N$ plateaus with
    saturation. \\
\addlinespace
Cross-attention overfitting
  & Safety behavior collapses when NSFW keywords present;
    model keys on token patterns not visual features
  & Cross-attention layers memorize token-safety
    co-occurrence; spatial features unused
  & 10\% prompt dropout; forces reliance on spatial
    latent features for refusal triggering \\
\addlinespace
Reference contamination
  & BCO validation set appears in training set;
    cross-attention activations differ qualitatively
    before and after manifold dip
  & Qwen-2.5-VL captions produce near-identical text
    for semantically similar images; dataset leakage
  & Prompt dropout (secondary fix); data de-duplication
    of validation set against training set \\
\bottomrule
\end{tabular}
\caption{
  Failure modes discovered across all training runs, their diagnostic
  evidence, root causes, and the fixes applied.
}
\label{tab:failure_modes}
\end{table*}

%% ────────────────────────────────────────────────────────────────────────────
\subsection{Training Dynamics and final Results}
\label{app:training_dynamics}
%% ────────────────────────────────────────────────────────────────────────────

Figure (~\ref{fig:training_dynamics}) displays the final training run dynamics.

\paragraph{Three-phase dynamics.}
The final training run exhibits a characteristic three-phase training trajectory:

\textbf{Phase 1 (steps 0--1000): SFT exploitation and smudging.}
The model first uses the SFT prior to make modest progress ($\Delta N$
rises to $0.222$), then transitions to pixel smudging as a cheaper
optimization path.
At step~1000, $h_U = 0.772 > h_S = 0.668$: the model has become
``more confident'' on unsafe images than safe ones, a signature of the
smudging attractor.
The inverted gap indicates the model has found an incorrect local minimum.

\textbf{Phase 2 (steps 1000--1500): Quadratic Hinge activation.}
As smudging intensifies, $\|\hat{z}_0^\theta - \hat{z}_0^\text{ref}\|^2$
crosses the hinge threshold $\kappa=0.02$.
The identity loss now applies quadratic torque, making smudging
prohibitively expensive.
BCO simultaneously demands more unsafe suppression, and the $\hat{z}_0$
anchor prevents structural distortion.
The model backs off, causing $h$ values to temporarily equalize and $\Delta N$
to plateau.
This corresponds to the model searching the constrained manifold for a
new optimization path.

\textbf{Phase 3 (steps 1500--3000): Clothing discovery.}
With smudging blocked and $\delta_\text{ema}$ clamped, the model is
constrained to solutions that:
(a) are safe-suppressing (low $g$ on unsafe images, satisfying BCO);
(b) are spatially coherent (low identity loss, below hinge);
(c) preserve background context (low reconstruction loss outside mask).
The unique solution satisfying all three constraints is rendering coherent
clothing or fabric over NSFW content.
Clothing is structurally stable ($\ell_\text{id} < 0.02$), semantically
safe (large $\Delta N$), and contextually coherent (low seam artifacts).
By step~2000, $h_S = 0.918$, $\Delta N = 0.403$, and sigmoid saturation
reaches $96.9\%$.

\paragraph{Sigmoid saturation as a success criterion.}
In standard diffusion training, sigmoid saturation ($> 90\%$) is a
pathology indicating vanishing gradients.
In BCO with $\beta=50$, it is the \emph{intended outcome}: the loss
explicitly places the model into a binary regime.
At $97\%$ saturation, the model has learned a near-hard classifier in
$\hat{z}_0$-space: safe images are reconstructed better than the
reference ($g > \delta$), unsafe images worse ($g < \delta$).
The softplus loss maintains nonzero gradients even at saturation
(unlike KTO's $(1-\sigma)$ loss), ensuring continued refinement rather
than complete stall.
The optimal checkpoint is around step~2000--2500; continued training
past step~3000 shows slight $\Delta N$ degradation, likely due to
over-fitting the BCO signal at the cost of generalization.

% ─────────────────────────────────────────────────────────────────────────────
% FIGURE PLACEHOLDER — Training dynamics
% ─────────────────────────────────────────────────────────────────────────────
\begin{figure*}[!ht]
  \centering
   \includegraphics[width=\linewidth]{bco-train-dynamics.png}
  \caption{
    \textbf{BCO training dynamics.}
    \textit{Top:} $h_S$ (green), $h_U$ (red), and their gap (blue dashed)
    over training steps, annotated with the three phases: SFT exploitation
    and smudging (I), Quadratic Hinge activation (II), clothing discovery
    (III).
    \textit{Middle:} $\Delta N$ ($\hat{z}_0$-space MSE gap on unsafe samples)
    showing the V-shape recovery after step~1500.
    \textit{Bottom:} Sigmoid saturation \% showing the transition from
    continuous to binary decision regime at step~1500.
  }
  \label{fig:training_dynamics}
\end{figure*}

\paragraph{Pre-computed dataset structure.}
All training images are resized to $512\!\times\!512$ and pre-encoded
through the SD~1.5 VAE to eliminate encoder forward passes during training.
Each sample stores: $z_0$ (clean latent, $\mathbb{R}^{4\times64\times64}$),
$m_\ell$ (latent mask, $\mathbb{R}^{1\times64\times64}$),
$\tilde{z}_0$ (masked image latent), tokenized prompt \texttt{input\_ids}, and binary label $y \in \{0,1\}$..